% slides.tex: a small slide layer for the TeX that runs on the engine.
%
% The format tex.wasm starts from has LaTeX 2020-02-02, the standalone class,
% xcolor and TikZ loaded, Computer Modern only, 75 fonts and 20,000 words of
% font metrics. There is no beamer, and TikZ costs 60 to 500 million engine
% steps per picture, so this file gives: fixed 160 x 90 mm pages, a slide
% environment, images and a few drawing commands as dvipdfmx specials.
%
% Font budget: every (font, size) pair is a font, and entering math mode at a
% size loads about ten. Keep `$', tabular and minipage[c] at \normalsize;
% use \sym for single symbols.
\scrollmode   % report errors in the log and carry on
\makeatletter
\hoffset=-1in \voffset=-1in
\newdimen\slidew \slidew=160mm
\newdimen\slideh \slideh=90mm
\textwidth=140mm \textheight=80mm
\newbox\slidebox \newbox\overbox \newbox\codebox
\newcount\slideno
\newif\iffootline \footlinetrue

% ------------------------------------------------------------------ look
% Colours are dvips colour specials written here directly: xcolor's
% \definecolor costs about ten million engine steps a colour in this format.
\def\defrgb#1#2{\expandafter\def\csname rgb@#1\endcsname{#2}}
\def\rgb#1{\csname rgb@#1\endcsname}
\defrgb{ink}{.102 .102 .102}    \defrgb{muted}{.604 .604 .604}
\defrgb{wasm}{.396 .310 .941}   \defrgb{accent}{.102 .051 .671}
\defrgb{alert}{.878 .224 .243}  \defrgb{alertdark}{.702 .179 .194}
\defrgb{codebg}{.961 .957 .941} \defrgb{boxbg}{.957 .953 .937}
\defrgb{boxhl}{.910 .941 .984}  \defrgb{boxhlb}{.290 .525 .910}
\defrgb{bar}{.184 .435 .839}    \defrgb{black20}{.8 .8 .8}
\defrgb{black60}{.4 .4 .4}      \defrgb{black70}{.3 .3 .3}
\defrgb{termbg}{.11 .12 .14}    \defrgb{termfg}{.80 .82 .84}
\defrgb{termcmd}{1 1 1}         \defrgb{termdim}{.50 .53 .56}
\defrgb{termhi}{.55 .85 .55}
\renewcommand\color[1]{\special{color push rgb \rgb{#1}}\aftergroup\rgb@pop\ignorespaces}
\def\rgb@pop{\special{color pop}}

% beamer's 11pt sizes, in Computer Modern Sans
\renewcommand\familydefault{\sfdefault}
\renewcommand\normalsize{\@setfontsize\normalsize{10.95}{13.6}}
\renewcommand\small{\@setfontsize\small{10}{12}}
\renewcommand\footnotesize{\@setfontsize\footnotesize{9}{11}}
\renewcommand\scriptsize{\@setfontsize\scriptsize{8}{9.5}}
\renewcommand\tiny{\@setfontsize\tiny{6}{7}}
\renewcommand\large{\@setfontsize\large{12}{14}}
\renewcommand\Large{\@setfontsize\Large{14.4}{18}}
\renewcommand\LARGE{\@setfontsize\LARGE{17.28}{22}}
\renewcommand\huge{\@setfontsize\huge{20.74}{25}}
\renewcommand\Huge{\@setfontsize\Huge{24.88}{30}}

\newcommand\muted[1]{{\color{muted}#1}}
\newcommand\link[1]{{\color{accent}#1}}
\newcommand\hi[1]{{\color{wasm}#1}}
\newcommand\red[1]{{\color{alertdark}#1}}
% one symbol from cmsy without entering math mode: \sym{"19} is \approx,
% \sym{"01} \cdot, \sym{"0F} \bullet
\newcommand\sym[1]{{\fontencoding{OMS}\fontfamily{cmsy}\fontseries{m}\fontshape{n}\selectfont\char#1}}
\def\${\char36 }   % the kernel's \$ wants a TS1 font, which is not there
\newcommand\strike[1]{\leavevmode\setbox\z@\hbox{#1}\rlap{\raise.55ex\hbox{\vrule width\wd\z@ height.07em}}\box\z@}
\renewcommand\labelitemi{\color{wasm}\raisebox{.15ex}{\small\sym{"0F}}}
\renewcommand\labelenumi{\color{wasm}\theenumi.}

% ------------------------------------------------------------------ frames
% \begin{slide}[Title] ... \end{slide}: one page. The body is centred
% vertically below the title; \vfill in the body overrides that.
\newcount\shipped
\def\titleskip{4mm}   % space above a slide's title
\newenvironment{slide}[1][]{%
  \global\advance\slideno\@ne \global\advance\shipped\@ne
  \setbox\slidebox\vbox to\slideh\bgroup
    \hsize\textwidth \linewidth\hsize \columnwidth\hsize
    \parindent\z@ \parskip\z@ \normalfont\normalsize
    \def\@tempa{#1}\ifx\@tempa\@empty\else
      \vspace*{\titleskip}{\centering\LARGE #1\par}\fi
    \vskip\z@ plus 1fil\relax
}{\par\vskip\z@ plus 1fil\egroup \shipslide}
% The standalone class ships pages of its own at \end{document}; after
% \nomorepages a \shipout builds its box and throws it away.
\def\nomorepages{\gdef\shipout{\setbox\z@}}
\def\skipslide{\global\advance\slideno\@ne}   % stands for a slide left out (pick.mjs)
\def\shipslide{\global\count0=\slideno   % TeX prints [page number] as it ships
  \shipout\vbox{\special{papersize=160mm,90mm}\color{ink}\offinterlineskip
    \vbox to\z@{\box\overbox \vss}%
    \iffootline\vbox to\z@{\vskip\dimexpr\slideh-5mm\relax
      \hbox to\slidew{\hfil\normalfont\scriptsize\color{muted}\the\slideno\hskip4mm}\vss}\fi
    \moveright10mm\box\slidebox}%
  \global\setbox\overbox\box\voidb@x}
\newcommand\statement[1]{\vfill\begin{center}\LARGE #1\end{center}\vfill}
% put something in the bottom right corner of the page (inside a slide)
\newcommand\atSE[1]{\global\setbox\overbox\vbox to\slideh{\vfil\hbox to\slidew{\hfil #1}}}
% a page that is one image, page wide
\newcommand\fullbleed[1]{\global\advance\slideno\@ne \global\advance\shipped\@ne
  \global\count0=\slideno
  \shipout\vbox{\special{papersize=160mm,90mm}\offinterlineskip
    \vbox to\slideh{\vss\hbox{\imgw{\slidew}{#1}}\vss}}}

% ------------------------------------------------------------------ images
% TeX never opens the image: figdims.tex (generated) gives the pixel size,
% and dvipdfmx embeds the file. \imgh{height}{file}, \imgw{width}{file};
% they set \imgwidth and \imgheight globally, for what is placed around it.
\def\figdim#1#2#3{\expandafter\gdef\csname fig@#1\endcsname{{#2}{#3}}}
\newdimen\imgwidth \newdimen\imgheight
\def\img@set#1#2{\def\img@W{#1}\def\img@H{#2}}
\def\img@dims#1{\@ifundefined{fig@#1}{\errmessage{no size for image #1}}{}%
  \expandafter\expandafter\expandafter\img@set\csname fig@#1\endcsname}
\newcommand\imgh[2]{\img@dims{#2}\global\imgheight=\dimexpr#1\relax
  \global\imgwidth=\dimexpr\imgheight*\img@W/\img@H\relax \img@put{#2}}
\newcommand\imgw[2]{\img@dims{#2}\global\imgwidth=\dimexpr#1\relax
  \global\imgheight=\dimexpr\imgwidth*\img@H/\img@W\relax \img@put{#2}}
\def\img@put#1{\leavevmode\hbox to\imgwidth{\vbox to\imgheight{\vfil
  \special{pdf:image width \the\imgwidth\space height \the\imgheight\space (#1)}}\hfil}}
% right after an image: put #3 at the point (#1,#2) of it, as fractions
\newcommand\imgmark[3]{\llap{\hbox to\imgwidth{\kern#1\imgwidth
  \raise#2\imgheight\hbox to\z@{#3\hss}\hss}}}

% ------------------------------------------------------------------ drawing
\def\bp#1{\strip@pt\dimexpr 0.99626\dimexpr#1\relax\relax}   % a length as a number of bp
% \ellipse{rx}{ry}{line width}: a red ellipse centred at the current point
\newcommand\ellipse[3]{\begingroup
  \dimen0=\dimexpr#1\relax \dimen2=\dimexpr#2\relax \dimen4=.55228\dimen0 \dimen6=.55228\dimen2
  \edef\e@a{\bp{\dimen0}}\edef\e@b{\bp{\dimen2}}\edef\e@ka{\bp{\dimen4}}\edef\e@kb{\bp{\dimen6}}%
  \special{pdf:content \rgb{alert} RG \bp{#3} w
    \e@a\space 0 m \e@a\space \e@kb\space \e@ka\space \e@b\space 0 \e@b\space c
    -\e@ka\space \e@b\space -\e@a\space \e@kb\space -\e@a\space 0 c
    -\e@a\space -\e@kb\space -\e@ka\space -\e@b\space 0 -\e@b\space c
    \e@ka\space -\e@b\space \e@a\space -\e@kb\space \e@a\space 0 c S}\endgroup}
% \tintbox{fill}{border}{rule width}{padding}{content}
\newcommand\tintbox[5]{\hbox{\setbox\z@\hbox{\kern#4\relax #5\kern#4}%
  \dimen@=\dimexpr\ht\z@+#4\relax \dimen@ii=\dimexpr\dp\z@+#4\relax
  \special{pdf:content \rgb{#1} rg \rgb{#2} RG \bp{#3} w 0 -\bp{\dimen@ii} \bp{\wd\z@}
    \bp{\dimexpr\dimen@+\dimen@ii\relax} re B}%
  \vrule width\z@ height\dimen@ depth\dimen@ii \box\z@}}
\newcommand\circled[3]{\leavevmode\setbox\z@\hbox{#3}\rlap{\kern.5\wd\z@\raise.3em\hbox{\ellipse{#1}{#2}{1pt}}}\box\z@}
\newcommand\rotated[2]{\leavevmode\special{pdf:btrans rotate #1}\hbox{#2}\special{pdf:etrans}}
% \flow[label]{length}: an arrow to the right along the baseline
\newcommand\flow[2][]{\hbox to#2{\rlap{\special{pdf:content \rgb{black60} RG \rgb{black60} rg .5 w
    0 0 m \bp{\dimexpr#2-1mm\relax} 0 l S \bp{#2} 0 m \bp{\dimexpr#2-2mm\relax} 1.6 l
    \bp{\dimexpr#2-2mm\relax} -1.6 l f}}\hss
  \raise2.5pt\hbox{\normalfont\tiny\ttfamily\color{black60}#1}\hss}}
% \stage{lines}: a box of centred \scriptsize lines (separated by \\),
% centred on the baseline; \stagehl is the highlighted one. \stagemark is
% drawn at the centre of the next \stagehl.
\let\stagemark\relax
\newcommand\stage[1]{\stage@{black20}{boxbg}{.4pt}{}{#1}}
\newcommand\stagehl[1]{\stage@{boxhlb}{boxhl}{.8pt}{\stagemark}{#1}}
\newcommand\stage@[5]{%
  \setbox\z@\vbox{\normalfont\scriptsize\let\\\cr\halign{\hfil##\hfil\cr#5\crcr}}%
  \dimen@=\wd\z@ \ifdim\dimen@<\dimexpr24mm-6pt\relax \dimen@=\dimexpr24mm-6pt\relax\fi
  \setbox\z@\hbox to\dimen@{\hss\vbox to\dimexpr12mm-6pt\relax{\vss\box\z@\vss}\hss}%
  \setbox\z@\hbox{\tintbox{#2}{#1}{#3}{3pt}{\box\z@}}%
  \dimen@=\dimexpr(\ht\z@-\dp\z@)/2\relax
  \lower\dimen@\hbox{\rlap{\kern.5\wd\z@\raise\dimen@\hbox{#4}}\box\z@}}
% \barrow{label}{width}{value}: one row of a bar chart
\newcommand\barrow[3]{\hbox{\hbox to9mm{\hss\bfseries #1}\kern2mm
  {\color{bar}\vrule height2.6mm depth1.4mm width#2}\kern1.5mm{\small\color{black70}#3}}}

% ------------------------------------------------------------------ code
% \begin{code}{font setup} ... \end{code}: lines as typed, on a tinted
% background. \ { } keep their meaning, so \hi{...} and \red{...} work.
\font\codenine=cmtt9
\font\codesmall=cmtt8 at 6.6pt
\font\codetiny=cmtt8 at 4.6pt
\def\code@sp{\leavevmode\kern\fontdimen2\font}
{\obeyspaces\gdef\code@spaces{\obeyspaces\let =\code@sp}}
\def\code#1{\par\setbox\codebox\vbox\bgroup #1\relax
  \hsize\dimexpr\linewidth-8pt\relax \parindent\z@ \parskip\z@
  \rightskip\z@ plus 1fil \parfillskip\z@ \frenchspacing
  \catcode`\%=12 \catcode`\_=12 \catcode`\~=12 \catcode`\$=12
  \catcode`\&=12 \catcode`\#=12 \catcode`\^=12
  \obeylines \code@spaces}
\def\endcode{\egroup \noindent\tintbox{codebg}{codebg}{0pt}{4pt}{\box\codebox}\par}
% \begin{term}{font setup} ... \end{term}: the same, as a dark terminal.
% \cmd{...} a command line, \note{...} a comment, \hl{...} a line to notice.
\def\term#1{\code{#1}\hsize\dimexpr\linewidth-12pt\relax}
\def\endterm{\egroup \noindent\tintbox{termbg}{termbg}{0pt}{6pt}{\color{termfg}\box\codebox}\par}
\newcommand\cmd[1]{{\color{termcmd}#1}}
\newcommand\note[1]{{\color{termdim}#1}}
\newcommand\hl[1]{{\color{termhi}#1}}
\makeatother
